\documentclass{article}

\usepackage{microtype}
\usepackage{graphicx}
\usepackage{booktabs}
\usepackage{amsmath}
\usepackage{algorithm}
\usepackage{algorithmic}
\usepackage{xcolor}
\usepackage{listings}
\usepackage{hyperref}

\providecommand{\theHalgorithm}{\arabic{algorithm}}
\usepackage[accepted]{icml2021}
\renewcommand{\ttdefault}{cmtt}

\icmltitlerunning{Contextual Bandit Event Recommendation for JoinIn}

\lstdefinelanguage{Dart}{
  keywords={class, enum, static, final, const, required, return, if, else, for, in, switch, case, try, catch, await, async, Future, List, Map, double, int, String, bool},
  keywordstyle=\color{blue!60!black}\bfseries,
  comment=[l]{//},
  commentstyle=\color{gray},
  stringstyle=\color{green!40!black},
  morestring=[b]',
}
\lstset{
  language=Dart,
  basicstyle=\ttfamily\tiny,
  breaklines=true,
  columns=fullflexible,
  frame=single,
  showstringspaces=false
}

\begin{document}

\twocolumn[
\icmltitle{Contextual Bandit Event Recommendation for JoinIn}

\begin{icmlauthorlist}
\icmlauthor{Marco Carraro}{unipd}
\icmlauthor{Francesco Vezzani}{unipd}
\end{icmlauthorlist}

\icmlaffiliation{unipd}{University of Padua, Padua, Italy}

\icmlcorrespondingauthor{Francesco Vezzani}{francesco.vezzani@studenti.unipd.it}

\icmlkeywords{Reinforcement Learning, Contextual Bandits, Event Recommendation, Flutter, Firebase}

\vskip 0.3in
]

\printAffiliationsAndNotice{}

\begin{abstract}
We build a reinforcement-learning event recommender for JoinIn, a Flutter/Firebase app where students find activities through a swipe feed. The existing Firestore query stays as a hard filter, and a Contextual Thompson Sampling re-ranker decides the order of the events that pass it. Ranking uses user, event, time, popularity, and safety features that the app already stores. User actions map to shaped rewards. For this class project we leave the remote logging hooks as no-ops and evaluate the policy offline with a synthetic simulator, instead of changing the Firebase backend.
\end{abstract}

\section{Goal and Setting}
JoinIn shows each user a short feed of candidate activities. Before this work, the feed was sorted mainly by event time, after dropping events that were expired, full, reported, owned by the user, or already joined. Our goal is to turn this fixed order into a learning recommender: one that favors events the user is likely to enjoy while still exploring uncertain ones.

Because real user interactions cannot be collected within the class-project window, we evaluate the recommender entirely on synthetic data. A local simulator generates users, events, and contextual signals that mirror the structure of the JoinIn feed, and a ground-truth utility function not exposed to the policy assigns actions to impressions using the same reward shaping the deployed app would log. This lets us exercise the full pipeline (feature extraction, ranking, reward mapping) without touching the Firebase backend or collecting new behavioral data from real students.

\section{RL Formulation}
At each step the app sees a user $u$ and a candidate event $e$ that passed the filter. We describe the pair with a context vector
\[
x_{u,e} = [\text{profile}(u), \text{event}(e), \text{time}(e), \text{social}(e), \text{safety}(e)].
\]
The action is the order of events in the feed. The reward comes from what the user does next with the card.

\begin{table}[t]
\caption{Reward shaping used by the implementation.}
\label{tab:rewards}
\vskip 0.1in
\begin{center}
\begin{small}
\begin{tabular}{lc}
\toprule
Action & Reward \\
\midrule
Impression & 0.0 \\
Join free activity & 1.0 \\
Request to join & 0.7 \\
Maybe/save & 0.3 \\
Pass/swipe left & -0.1 \\
Leave after join & -0.6 \\
Report activity & -1.0 \\
Chat activity after join & 0.2 \\
\bottomrule
\end{tabular}
\end{small}
\end{center}
\vskip -0.1in
\end{table}

\section{Algorithm Choice}
We use Contextual Thompson Sampling with a linear reward model \cite{agrawal2013thompson,li2010contextual}:
\[
\hat{r}(u,e) = x_{u,e}^{T}\theta.
\]
For each candidate, the app draws a perturbed weight vector $\tilde{\theta}$ from a light posterior approximation and ranks by $x_{u,e}^{T}\tilde{\theta}$. High-scoring events go first, but uncertain yet plausible events still get a chance to appear.

We chose this method because JoinIn has a changing event catalog, sparse feedback, and strict product rules. A contextual bandit fits well: each swipe is a single one-step decision, the candidate set turns over quickly, and a first deployment must stay interpretable and safe. The generic ranking step comes from the published Dart package \texttt{dart\_rl} version 1.0.0 \cite{dartRlPackage}; the app only adds JoinIn-specific features, priors, and logging.

\begin{algorithm}[tb]
\caption{Feed re-ranking with Contextual Thompson Sampling}
\label{alg:cts}
\begin{algorithmic}
\STATE {\bfseries Input:} user $u$, filtered events $E$
\FOR{each event $e \in E$}
\STATE Build feature vector $x_{u,e}$
\STATE Sample weights $\tilde{\theta}$ around the current prior
\STATE Compute score $s_e = x_{u,e}^{T}\tilde{\theta}$
\ENDFOR
\STATE Sort events by $s_e$ descending
\STATE Log one impression document per shown event
\STATE On user action, log action and shaped reward
\end{algorithmic}
\end{algorithm}

\section{Implementation}
\subsection{Codebase Integration}
We take the reusable linear Thompson Sampling code from the published \texttt{dart\_rl} package and keep all product-specific logic in JoinIn:
\begin{itemize}
    \item \textbf{Dependency}: \texttt{pubspec.yaml} adds \texttt{dart\_rl: \^{}1.0.0}.
    \item \textbf{Reusable package}: \texttt{dart\_rl} provides the linear Thompson Sampling sampler, bandit-arm wrapper, and ranking helper used by the service.
    \item \textbf{Recommendation service}: feature extraction, prior weights, reward mapping, and logging hooks for the local simulator.
    \item \textbf{Home feed}: runs the re-ranker after candidate retrieval and records feed interactions.
\end{itemize}

Candidate retrieval stays in \texttt{FirebaseApi.fetchEventsForFeed}. The query filters by future date, visible universities, ownership, participation, report threshold, user reports, and capacity. The re-ranker only ever sees candidates that already passed these checks.

\subsection{State Features}
The feature vector is small and uses only data already in \texttt{EventModel} and \texttt{UserModel}: a bias term, university match, event type, time-to-event, evening and weekend flags, attendee ratio, waiting-list ratio, image availability, freshness, profile-category match, and a report penalty.

\begin{lstlisting}
static Map<String, double> buildFeatures({
  required UserModel user,
  required EventModel event,
}) {
  final attendeeCount = event.attendees?['uids']?.length ?? 0;
  final waitingCount = event.waitings?['uids']?.length ?? 0;
  final maxAttendees = event.maxAttendees?.toDouble();

  return {
    'bias': 1.0,
    'same_university': event.university == user.university ? 1.0 : 0.0,
    'free_to_join': event.type == EventType.freeToJoin.toString() ? 1.0 : 0.0,
    'request_to_join': event.type == EventType.requestToJoin.toString() ? 1.0 : 0.0,
    'attendee_ratio': _clamp01(attendeeCount / capacity),
    'waiting_ratio': _clamp01(waitingCount / capacity),
    'report_penalty': _clamp01((event.reportCount ?? 0) / 5.0),
  };
}
\end{lstlisting}

\subsection{Ranking}
The service holds an interpretable prior weight vector and one standard deviation per feature. On each feed refresh it hands the ranking step to \texttt{dart\_rl}, which samples around the prior and sorts events by the sampled score.

\begin{lstlisting}
static List<RankedEvent> rankEvents({
  required List<EventModel> events,
  required UserModel user,
}) {
  final bandit = LinearThompsonSampling<EventModel>(
    weights: _priorWeights,
    standardDeviations: _posteriorStd,
  );
  final arms = events.map((event) => ContextualBanditArm(
    item: event,
    features: buildFeatures(user: user, event: event),
  )).toList();
  return bandit.rank(arms).map((ranking) {
    return RankedEvent(
      event: ranking.item,
      score: ranking.score,
      exploitationScore: ranking.expectedReward,
      explorationBonus: ranking.explorationBonus,
      features: ranking.features,
    );
  }).toList();
}
\end{lstlisting}

\subsection{No Server-Side Logging}
Our early analysis showed that real pass and maybe signals live only in the local Isar store, and that production training would need centralized logs. For this class project we chose not to touch the Firebase backend or add new server collections. So the app-side logging methods are no-ops by default, and we evaluate with the local synthetic simulator included in the submission folder.

\begin{lstlisting}
static Future<void> logInteraction({
  required UserModel? user,
  required EventModel event,
  required RecommendationAction action,
  required String sessionId,
}) async {
  debugPrint('Skipped remote logging; synthetic simulations are used.');
}
\end{lstlisting}

\section{Evaluation Plan and Results}
We checked that the implementation builds: \texttt{flutter pub get} resolves the published \texttt{dart\_rl} 1.0.0 package from pub.dev, and the Flutter analyzer reports only pre-existing style hints in \texttt{home\_page.dart}, not type errors from the recommender.

Because we do not change the Firebase server, evaluation runs through the local script \texttt{simulations/synthetic\_bandit\_simulation.py}. Training on real interactions would mean backend changes and a privacy-policy update, since we would collect and process new behavioral signals for personalization. That review takes longer than the class-project window, so we use synthetic data now and plan a privacy-reviewed release for Q4. The simulator creates 1{,}200 users with 30 candidate events each. Each policy shows the top five events per user, giving 6{,}000 impressions per policy. Users have a university and a preferred category; events have category, university, type, time, popularity, freshness, image availability, and report risk. A ground-truth utility function not exposed to the policy then picks each action using the same reward shaping as the app.

The baseline sorts candidates by earliest event time. The proposed method ranks the same candidates with Contextual Thompson Sampling. The run writes \texttt{results/synthetic\_metrics.json} and \texttt{results/synthetic\_interactions.csv}.

\begin{table}[t]
\caption{Synthetic simulation results. Each policy is evaluated on 6,000 impressions. Delta is CTS minus baseline.}
\label{tab:results}
\vskip 0.1in
\begin{center}
\begin{small}
\begin{tabular}{lccc}
\toprule
Metric & Baseline & CTS & Delta \\
\midrule
Join rate & 40.97\% & 58.10\% & +17.13 pp \\
Request rate & 15.43\% & 10.95\% & -4.48 pp \\
Pass rate & 43.60\% & 30.95\% & -12.65 pp \\
Avg. reward/impression & 0.474 & 0.627 & +0.153 \\
\bottomrule
\end{tabular}
\end{small}
\end{center}
\vskip -0.1in
\end{table}

The bandit policy raises the join rate and the average reward and lowers the pass rate. The request rate drops because the learned ranking favors free-to-join events, which carry the strongest positive reward in the synthetic user model. That is fine for a first version, since a direct join is the best outcome, but a production system should watch that request-to-join events do not become under-exposed.

\section{Limitations and Next Steps}
The current model uses a fixed prior and on-device random exploration. This is safe for a first deployment and fits a class project that avoids server changes, but it does not yet update the posterior from real user rewards. The synthetic simulator checks the mechanics and the reward design; it cannot prove a production lift. The Q4 release should: update the privacy policy and backend data contract; add an opt-in analytics pipeline; estimate $\theta$ from real past interactions; and ship updated weights through Remote Config or a server endpoint. It should also model slate effects and position bias, since an event's reward depends on where it sits in the stack.

\bibliography{main}
\bibliographystyle{icml2021}

\end{document}
